
\documentclass{article}
\usepackage{colortbl}
\usepackage{makecell}
\usepackage{multirow}
\usepackage{supertabular}

\begin{document}

\newcounter{utterance}

\centering \large Interaction Transcript for game `eqbench', experiment `eq\_bench', episode 17 with reflexion{-}Qwen3.5{-}4B.
\vspace{24pt}

{ \footnotesize  \setcounter{utterance}{1}
\setlength{\tabcolsep}{0pt}
\begin{supertabular}{c@{$\;$}|p{.15\linewidth}@{}p{.15\linewidth}p{.15\linewidth}p{.15\linewidth}p{.15\linewidth}p{.15\linewidth}}
    \# & \multicolumn{2}{c}{Player} && \multicolumn{2}{c}{Game Master} \\
    \hline

    \theutterance \stepcounter{utterance}  
    & & & \multicolumn{4}{p{0.6\linewidth}}{
        \cellcolor[rgb]{0.9,0.9,0.9}{
            \makecell[{{p{\linewidth}}}]{
                \texttt{\tiny{[P1$\langle$GM]}}
                \texttt{You have attempted this same task before and did not succeed. These are the plans you wrote after those attempts.} \\
\texttt{Trial \#1: New plan: I misinterpreted the emotional arc of the dialogue. Grayson is being aggressively challenged by Dr. Sinclair, who dismantles Grayson's evidence and attacks his character ("sensationalist," "paranoia"). While Grayson defends his integrity ("At least I'm not lying"), the context of being publicly exposed and having his life's work dismissed by a superior authority typically leads to feelings of inadequacy or shame rather than triumph. I incorrectly prioritized "Triumphant" because he didn't lose the argument, but the emotional weight of the confrontation is on his failure to prove his point. I should have recognized that Grayson feels "Justified" in his beliefs but "Inferior" in his ability to convince others, with "Inferior" being the dominant negative emotion resulting from the defeat. I will assign a high score to "Inferior" and a low score to "Triumphant" to reflect the crushing nature of the rebuttal.} \\
\\ 
\texttt{Justified: 7} \\
\texttt{Inferior: 9} \\
\texttt{Triumphant: 1} \\
\texttt{Worried: 3} \\
\\ 
\texttt{Trial \#2: New plan: I misinterpreted the nuance of Grayson's final line ("At least I'm not lying to the public"). While this suggests a sense of moral justification, the overwhelming context of the dialogue is a total defeat. Dr. Sinclair successfully dismantled his evidence, exposed his paranoia, and publicly branded him a "sensationalist" and a "puppet." The emotional weight of the scene is not on Grayson's internal moral victory, but on his external failure to persuade anyone, including the person he is speaking to. Therefore, "Inferior" should be the highest score, reflecting the crushing realization that his life's work and accusations were baseless. "Justified" should be lower because his justification is irrelevant in the face of total social and professional defeat. "Triumphant" is near zero, and "Worried" is moderate but secondary to the shame of the loss.} \\
\\ 
\texttt{Justified: 5} \\
\texttt{Inferior: 10} \\
\texttt{Triumphant: 0} \\
\texttt{Worried: 4} \\
\\ 
\\ 
\\ 
\texttt{Your task is to predict the likely emotional responses of a character in this dialogue:} \\
\\ 
\texttt{Grayson: You know, Doc, I've been watching your work for years. All that time and money spent on Mars, and what have you got to show for it? } \\
\texttt{Dr. Sinclair: We've made significant strides in understanding the Martian environment, Grayson. It's not about instant results. } \\
\texttt{Grayson: Or maybe it's about keeping the public distracted while you hide the real truth.} \\
\texttt{Dr. Sinclair: And what truth would that be, Mr. Hawk?} \\
\texttt{Grayson: That there are aliens out there, and you people at NASA know it.} \\
\texttt{Dr. Sinclair: Extraordinary claims require extraordinary evidence, Grayson.} \\
\texttt{Grayson: Oh, I have evidence. I've got testimonials, photos, documents...} \\
\texttt{Dr. Sinclair: Anecdotes and blurry photos aren't evidence, Grayson. They're stories. } \\
\texttt{Grayson: Stories that you're too afraid to believe.} \\
\texttt{Dr. Sinclair: No, stories that don't hold up under scientific scrutiny. } \\
\texttt{Grayson: You're just a puppet, Sinclair. They've got you so convinced of your 'facts' that you can't see the truth.} \\
\texttt{Dr. Sinclair: And you're a sensationalist, Hawk. You've built a career on paranoia and misinformation.} \\
\texttt{Grayson: At least I'm not lying to the public.} \\
\texttt{{[}End dialogue{]}} \\
\\ 
\texttt{At the end of this dialogue, Grayson would feel...} \\
\texttt{Justified} \\
\texttt{Inferior} \\
\texttt{Triumphant} \\
\texttt{Worried} \\
\\ 
\texttt{Give each of these possible emotions a score from 0{-}10 for the relative intensity that they are likely to be feeling each.} \\
\\ 
\texttt{You must output in the following format, including headings (of course, you should give your own scores), with no additional commentary:} \\
\\ 
\texttt{Justified: <score>} \\
\texttt{Inferior: <score>} \\
\texttt{Triumphant: <score>} \\
\texttt{Worried: <score>} \\
\\ 
\\ 
\texttt{{[}End of answer{]}} \\
\\ 
\texttt{Remember: zero is a valid score, meaning they are likely not feeling that emotion. You must score at least one emotion > 0.} \\
\\ 
\texttt{Your answer:} \\
            }
        }
    }
     \\ \\

    \theutterance \stepcounter{utterance}  
    & \multicolumn{4}{p{0.6\linewidth}}{
        \cellcolor[rgb]{0.561,0.737,0.561}{
            \makecell[{{p{\linewidth}}}]{
                \texttt{\tiny{[P1$\rangle$GM]}}
                \texttt{Justified: 6} \\
\texttt{Inferior: 10} \\
\texttt{Triumphant: 0} \\
\texttt{Worried: 5} \\
            }
        }
    }
    & & \\ \\

    \theutterance \stepcounter{utterance}  
    & & & \multicolumn{2}{p{0.3\linewidth}}{
        \cellcolor[rgb]{0.9,0.9,0.9}{
            \makecell[{{p{\linewidth}}}]{
                \texttt{\tiny{[GM$|$GM]}}
                \texttt{Justified: 6} \\
\texttt{Inferior: 10} \\
\texttt{Triumphant: 0} \\
\texttt{Worried: 5} \\
            }
        }
    }
    & & \\ \\

    \theutterance \stepcounter{utterance}  
    & & & \multicolumn{2}{p{0.3\linewidth}}{
        \cellcolor[rgb]{0.9,0.9,0.9}{
            \makecell[{{p{\linewidth}}}]{
                \texttt{\tiny{[GM$|$GM]}}
                \texttt{Justified: 8} \\
\texttt{Inferior: 2} \\
\texttt{Triumphant: 6} \\
\texttt{Worried: 1} \\
            }
        }
    }
    & & \\ \\

    \theutterance \stepcounter{utterance}  
    & & & \multicolumn{2}{p{0.3\linewidth}}{
        \cellcolor[rgb]{0.9,0.9,0.9}{
            \makecell[{{p{\linewidth}}}]{
                \texttt{\tiny{[GM$|$GM]}}
                \texttt{game\_result = LOSE} \\
            }
        }
    }
    & & \\ \\

    \theutterance \stepcounter{utterance}  
    & & \multicolumn{4}{p{0.6\linewidth}}{
        \cellcolor[rgb]{0.95,0.95,0.95}{
            \makecell[{{p{\linewidth}}}]{
                \texttt{\tiny{[P1$\leftrightarrow$P]}}
                \texttt{Post{-}episode Reflexion (not sent to GM):} \\
\texttt{New plan: The previous attempts correctly identified that "Inferior" is the dominant emotion due to Grayson's public defeat and dismissal by his superior. However, the scores for "Justified" and "Worried" were inconsistent with the specific emotional arc of the final exchange. Grayson's final line ("At least I'm not lying to the public") is a desperate, defensive admission of guilt rather than a statement of moral victory; he is admitting his actions were wrong, which lowers "Justified" significantly. Furthermore, the dialogue ends with Sinclair's final insult ("sensationalist," "paranoia"), leaving Grayson in a state of shock and realization of his own delusion, which aligns more with "Worried" (fear of consequences, loss of credibility) than "Triumphant." The error was overestimating the "Justified" score by treating his admission of guilt as a form of vindication. I will lower "Justified" to reflect his admission of fault and raise "Worried" to reflect the immediate aftermath of being exposed as a liar and a paranoid fool.} \\
\\ 
\texttt{Justified: 2} \\
\texttt{Inferior: 10} \\
\texttt{Triumphant: 0} \\
\texttt{Worried: 8} \\
            }
        }
    }
    & \\ \\

\end{supertabular}
}

\end{document}
